\clearpage
% C057 | 考前 30 秒复习卡
\begin{summarybox}
  \textbf{\textcolor{CSummary}{一句话背法}}

  \textbf{过焦点，先看支：同支倒数加，异支倒数减；都等于 $2a/b^2$。}
  \[
    \boxed{\begin{array}{ll}
       \text{同支：}&\displaystyle\frac1{|AF|}+\frac1{|BF|}=\frac{2a}{b^2},\\[7pt]
       \text{异支：}&\displaystyle\left|\frac1{|AF|}-\frac1{|BF|}\right|=\frac{2a}{b^2}.
    \end{array}}
  \]

  \medskip
  \textbf{\textcolor{CSummary}{三大识别信号}}
  \begin{enumerate}[leftmargin=*,itemsep=0.35em,topsep=0.3em]
    \item 标准双曲线 $x^2/a^2-y^2/b^2=1$，$a,b>0$，$c^2=a^2+b^2$。
    \item 直线过一个焦点 $F(\pm c,0)$，\textbf{有两个不同的有限交点} $A,B$。
    \item 同支时 $F$ 在 $AB$ 内；异支时 $F$ 在 $AB$ 的\textbf{延长线}上。
  \end{enumerate}

  \medskip
  \textbf{\textcolor{CSummary}{想不起证明？只抓两个式子}}

  令 $r=|PF|$、$e=c/a$、$p=b^2/a$。右焦点对应
  \[
    \text{右支 }r=ex-a,\qquad \text{左支 }r=a-ex.
  \]
  因为同一直线上的点满足\textbf{横向位移/焦距的比值同向相等、反向相反}，所以这两个式子相减或相加后，角度信息消失，只剩 $2/p=2a/b^2$。

  \medskip

  \textbf{\textcolor{CSummary}{最易忘的一条：焦准成对，取点自由。}}

  左焦点 $F_1(-c,0)$ 配左准线 $\ell_1:x=-a^2/c$；
  右焦点 $F_2(c,0)$ 配右准线 $\ell_2:x=a^2/c$。\textbf{点在任意一支，两套配对都成立}：
  \[
    \boxed{\frac{|PF_1|}{d(P,\ell_1)}
       =\frac{|PF_2|}{d(P,\ell_2)}
       =e=\frac ca>1}.
  \]
  \textbf{口诀：}“\textbf{焦准左右配，曲线取点不分支；选定一套不换边。}”
  异支的 $A,B$ 若直线经过右焦点，两点都可以用\textbf{右焦点--右准线}。
  注意比值是 $e$，不是 $1$，也不能左右交叉配对。

  \medskip
  \textbf{\textcolor{CSummary}{防错五句}}
  \begin{itemize}[leftmargin=*,itemsep=0.3em,topsep=0.3em]
    \item 焦准配对与点无关：\textbf{右焦点--右准线、左焦点--左准线}；点可任意支，距离比为 $e>1$。
    \item 右焦点异支时：\textbf{右支点更近}，$1/r_{\text{右}}-1/r_{\text{左}}=2a/b^2$。
    \item 同支已知一段 $r$：一定有 $r>b^2/(2a)$，否则不存在另一\textbf{有限}同支交点。
    \item $k=\pm b/a$ 为渐近线方向，只有一个有限交点，\textbf{不套本定理}。
    \item 做题\textbf{先算 $K=2a/b^2$}，再解未知段；不必先联立双曲线。
  \end{itemize}
\end{summarybox}
